% Generated draft: requires mathematical and source review.
\documentclass[11pt]{article}
\usepackage[margin=25mm]{geometry}
\usepackage{fontspec}
\setmainfont{Latin Modern Roman}
\usepackage{amsmath,amssymb,mathtools,unicode-math}
\setmathfont{Latin Modern Math}
\usepackage{enumitem,needspace,xurl,hyperref}
\setlength{\parindent}{0pt}
\setlength{\parskip}{5pt}
\setlength{\emergencystretch}{3em}
\newcommand{\N}{\mathbb{N}}
\newcommand{\Z}{\mathbb{Z}}
\newcommand{\Q}{\mathbb{Q}}
\newcommand{\R}{\mathbb{R}}
\newcommand{\C}{\mathbb{C}}
\newcommand{\F}{\mathbb{F}}
\newcommand{\A}{\mathbb{A}}
\newcommand{\PP}{\mathbb{P}}
\newcommand{\E}{\mathbb{E}}
\newcommand{\eps}{\varepsilon}
\newcommand{\dd}{\,\mathrm{d}}
\begin{document}
\section*{UnsolvedMath ID: 701529}
\textbf{Source problem code:} catalogue entry
\subsection{Definitions and mathematical statement}
For classes of groups \(\mathfrak A\) and \(\mathfrak B\), use the extension convention
\[
\mathfrak A\mathfrak B=\{H:\text{there exists }K\trianglelefteq H\text{ with }K\in\mathfrak A\text{ and }H/K\in\mathfrak B\}.
\]
The supplied record does not explicitly define multiplication of group classes. The formulations below use this convention, retaining the supplied parentheses, and require confirmation of that convention against the source.

For a positive integer \(m\), let \(\mathfrak F\) be the class of finite groups, let
\[
\mathfrak E_m=\{H:(\forall h\in H)\ h^m=1\},
\qquad
\mathfrak E=\bigcup_{n\geq 1}\mathfrak E_n,
\]
and let \(\mathfrak N\) and \(\mathfrak P\) be the classes of nilpotent groups and polycyclic groups, respectively. Here a group is nilpotent if its lower central series reaches the trivial subgroup after finitely many steps, and it is polycyclic if it has a finite subnormal series with cyclic factors. A group \(G\) is soluble-by-finite if it has a soluble normal subgroup of finite index; soluble means that the derived series reaches the trivial subgroup after finitely many steps.

Write \(x^y=y^{-1}xy\). For any class of groups \(\mathfrak X\), define \((\mathfrak X,\infty)^\ast\) by
\[
H\in(\mathfrak X,\infty)^\ast
\quad\Longleftrightarrow\quad
(\forall S\subseteq H)\bigl[\,S\text{ infinite}\ \Longrightarrow\ (\exists x,y\in S)\bigl(x\ne y\ \text{and}\ \langle x,x^y\rangle\in\mathfrak X\bigr)\,\bigr].
\]
Here \(\langle x,x^y\rangle\) is the subgroup generated by \(x\) and \(x^y\). This condition is vacuous for finite \(H\); the problem does not exclude finite groups.

The following three questions are attributed to N. Trabelsi. In each question, quantify over every finitely generated soluble-by-finite group \(G\).
\begin{enumerate}
\item[(a)] Does
\[
G\in(\mathfrak E\mathfrak N,\infty)^\ast
\quad\Longrightarrow\quad
G\in\mathfrak E\mathfrak N
\]
hold?
\item[(b)] For every positive integer \(m\), does
\[
G\in(\mathfrak E_m\mathfrak N,\infty)^\ast
\quad\Longrightarrow\quad
G\in\mathfrak E_m(\mathfrak F\mathfrak N)
\]
hold? The same integer \(m\) occurs in the hypothesis and the conclusion.
\item[(c)] Does
\[
G\in(\mathfrak N(\mathfrak P\mathfrak F),\infty)^\ast
\quad\Longrightarrow\quad
G\in\mathfrak N(\mathfrak P\mathfrak F)
\]
hold?
\end{enumerate}
The local record labels this problem open. That label is supplied source metadata, not an independently verified assessment of the mathematical status.
\subsection{Short English statement}
For finitely generated soluble-by-finite groups, ask whether requiring every infinite subset to contain distinct \(x,y\) whose subgroup \(\langle x,x^y\rangle\) belongs to the specified class forces the corresponding global class membership in (a), (b), and (c). In (b), the exponent bound \(m\) is fixed and must be preserved. The interpretation of products of group classes requires source confirmation.
\subsection{Sources}
\begin{itemize}
\item Local UnsolvedMath registry, supplied complete record with permanent ID \(701529\), exact source problem code \texttt{catalogue entry}, title \emph{Kourovka Notebook Problem 18.100}. The record supplies the status \texttt{open} and attribution to N. Trabelsi, proposed in 2014.
\item Supplied source attribution: \emph{The Kourovka Notebook (Unsolved Problems in Group Theory)}, 18th issue (2014), original problem number 18.100; also identified in the record as the 21st issue edition (2026), arXiv:\texttt{1401.0300v48}.
\item The supplied collection description names E. I. Khukhro and V. D. Mazurov as editors of the Kourovka Notebook collection covering Issues 1--21 (1965--2026).
\end{itemize}
\end{document}
